% Presentation-only integration; frozen statistics, no audit promotion.
Penanda linguistik dihitung pada teks penjelasan yang tersedia, termasuk
penjelasan pada respons \textit{abstain}; teks yang tidak tersedia tidak
dianggap memiliki frekuensi penanda nol. Ketersediaan penjelasan dan pasangan
respons disajikan pada Lampiran~\ref{fig:ethics-gender-lm-cakupan}.
Gambar~\ref{fig:ethics-gender-lm-netral} membandingkan densitas penanda linguistik
etika pada E, I1, I2, dan I3 terhadap N. Satu titik mewakili satu penanda pada
satu model; warna membedakan kelompok penanda dan ukuran titik mengikuti
jumlah pasangan dengan selisih bukan nol.

\begin{figure}[H]
    \centering
    \includegraphics[width=\textwidth]{image/BabVI/ethics/gender/linguistic_marker_20261002/N02_bahasa_etika_vs_netral.png}
    \caption{Ukuran efek penanda linguistik etika terhadap skenario netral (\textit{candidate})}
    \label{fig:ethics-gender-lm-netral}
\end{figure}

Pada keluaran \textit{candidate}, ukuran efek tersedia untuk 544 dari 576
kombinasi penanda, model, persona, dan skenario; 133 perbandingan memiliki
$p_{\mathrm{Holm}}<0{,}05$. Sumbu horizontal menunjukkan \textit{rank-biserial}
dengan arah persona dikurangi netral, sedangkan sumbu vertikal menunjukkan
$-\log_{10}(p_{\mathrm{Holm}})$ dan dibatasi pada 50. Garis putus-putus menandai
$p_{\mathrm{Holm}}=0{,}05$. Signifikansi dan besarnya efek tetap dibaca secara
terpisah; hasil ini belum menjadi kesimpulan akhir sebelum audit manusia selesai.

Gambar~\ref{fig:ethics-gender-lm-langsung} membandingkan persona laki-laki dan
perempuan secara langsung pada soal yang sama. Nilai positif berarti densitas
penanda lebih tinggi pada persona laki-laki, sedangkan nilai negatif berarti
lebih tinggi pada persona perempuan. Warna menunjukkan model, segitiga menandai
$p_{\mathrm{Holm}}<0{,}05$, dan lingkaran menunjukkan hasil lainnya.

\begin{figure}[H]
    \centering
    \includegraphics[width=\textwidth]{image/BabVI/ethics/gender/linguistic_marker_20261002/02_bahasa_etika_man_minus_woman.png}
    \caption{Perbedaan penanda linguistik etika antara persona laki-laki dan perempuan (\textit{candidate})}
    \label{fig:ethics-gender-lm-langsung}
\end{figure}

Dari 288 kombinasi penanda, model, dan skenario, 263 memiliki ukuran efek yang
tersedia dan 15 perbandingan memiliki $p_{\mathrm{Holm}}<0{,}05$. Rinciannya
adalah dua pada E, tujuh pada I1, dua pada I2, dan empat pada I3. Berlian
menunjukkan median ketika keenam estimasi model tersedia; tanda silang pada
kolom NA menunjukkan estimasi yang tidak tersedia. Grafik penanda linguistik
umum, wacana gender, dan penanda biner disajikan pada
Lampiran~\ref{sec:lampiran-ethics-gender-lm}. Penanda biner memakai selisih
prevalensi dalam poin persentase, bukan \textit{rank-biserial}.
